%! AP Statistics — Unit 1, Lesson 1.4 — Group Activity ANSWER KEY
\documentclass[10pt]{article}
\usepackage{apstats-article}
\usepackage{apstats-key}

%=============================================================================
\begin{document}

\pageheader{Unit 1, Lesson 1.4}{Group Activity --- Answer Key}
\begin{tabularx}{\linewidth}{X X X}
Name:\;\ans{Key} & Partner:\; & Period:\;
\end{tabularx}

\vspace{0.1in}

% ── PART 1 ───────────────────────────────────────────────────────────────────
{\large\bfseries\color{navy} Part 1 \quad Construct the Bar Chart}

\vspace{0.1in}

\small\noindent Streaming service survey, $n = 50$.

\vspace{0.1in}

\renewcommand{\arraystretch}{2.0}
\begin{tabularx}{\linewidth}{@{} >{\bfseries}p{0.22\linewidth}
                                    C{0.18\linewidth}
                                    C{0.22\linewidth}
                                    C{0.22\linewidth} @{}}
\rowcolor{sky}
\textbf{Service} & \textbf{Frequency} & \textbf{Rel.\ Frequency} & \textbf{Angle (${}^\circ$)} \\
\midrule
Netflix  & 21 & \ans{0.42 (42\%)} & \ans{151.2} \\
Disney+  & 10 & \ans{0.20 (20\%)} & \ans{72}    \\
Hulu     & 9  & \ans{0.18 (18\%)} & \ans{64.8}  \\
HBO Max  & 6  & \ans{0.12 (12\%)} & \ans{43.2}  \\
Other    & 4  & \ans{0.08 (8\%)}  & \ans{28.8}  \\
\midrule
\rowcolor{sky}\textbf{Total} & \ans{50} & \ans{1.00} & \ans{360} \\
\end{tabularx}

\begin{teachernote}
Angles: Netflix $= 0.42 \times 360 = 151.2^\circ$; Disney+ $= 72^\circ$;
Hulu $= 64.8^\circ$; HBO Max $= 43.2^\circ$; Other $= 28.8^\circ$.
Sum $= 360^\circ$. Accept reasonable rounding.
\end{teachernote}

\vspace{0.1in}

\begin{teachernote}
\textbf{Correct bar chart:} 5 separated bars labeled Netflix (tallest, 42\%),
Disney+ (20\%), Hulu (18\%), HBO Max (12\%), Other (shortest, 8\%).
Y-axis: ``Relative Frequency'' or ``Proportion'', scale 0 to at least 0.45
(or 0 to 45\%). Bars do NOT touch. Title present (e.g., ``Preferred Streaming
Service for 50 AP Statistics Students'').

\textbf{Common errors:} bars touching; y-axis not starting at 0; missing title
or axis labels; bars arranged in an order other than the table (accept any order,
but consistent is best).
\end{teachernote}

\vspace{0.1in}

\begin{enumerate}[leftmargin=1.5em, itemsep=8pt]
  \item Sum $= 1$? \ans{Yes (allow small rounding error, e.g., 0.999 or 1.001).}

  \item \ans{Of the 50 AP Statistics students surveyed, Netflix was by far the
        most popular streaming service (42\%), while Other was the least popular
        (8\%).}

  \begin{teachernote}
  Require: $n = 50$, variable in context (streaming service), most and least
  common with relative frequencies. Deduct for raw counts only.
  \end{teachernote}

  \item \ans{The distribution is skewed toward Netflix, with no other service
        reaching even half of Netflix's share. Disney+ and Hulu are fairly close
        at 20\% and 18\%, while HBO Max and Other together account for only 20\%.}

\end{enumerate}

\vspace{0.12in}

% ── PART 2 ───────────────────────────────────────────────────────────────────
{\large\bfseries\color{navy} Part 2 \quad Misleading Graphs}

\vspace{0.1in}

\begin{enumerate}[leftmargin=1.5em, itemsep=8pt]

  \item \ans{Starting the y-axis at 30\% (0.30) instead of 0 means the bars for
        Disney+ (20\%) and Hulu (18\%) would not even appear --- they fall below
        the axis cutoff. Netflix's bar would look dramatically taller than the
        others. The graph would mislead a reader into thinking Netflix is used
        almost exclusively, when in fact Disney+, Hulu, and others account for
        58\% of responses combined.}

  \begin{teachernote}
  Students may describe a slightly different scenario (e.g., starting at 5\% or
  10\%). The key reasoning is the same: the y-axis cutoff makes the differences
  look larger than they are. Accept any reasonable specific example.
  \end{teachernote}

  \item \ans{A 3-D pie chart tilts the Netflix slice toward the viewer, making
        it appear even larger than its true 42\% share. Readers judge area or
        angle visually, and perspective distortion causes them to overestimate
        the size of near slices. The underlying data are correct, but the
        presentation inflates the apparent popularity of Netflix.}

  \item \ans{Bar chart: Netflix (42\%) is clearly the largest, followed by
        Disney+ (20\%), Hulu (18\%), HBO Max (12\%), Other (8\%). With 5
        categories, the bar chart is easier to compare than a pie chart because
        readers can compare bar heights directly rather than estimating sector
        angles. With only 5 categories, either chart is acceptable, but bar
        charts are generally preferred for comparison tasks.}

\end{enumerate}

\vspace{0.12in}

% ── PART 3 ───────────────────────────────────────────────────────────────────
{\large\bfseries\color{navy} Part 3 \quad Reflection}

\vspace{0.1in}

\begin{tcolorbox}[colback=goldbg, colframe=goldacc, boxrule=0.8pt, arc=2mm,
    left=3mm, right=3mm, top=2mm, bottom=2mm]
\small
\ans{The most important thing to check is whether the y-axis starts at 0.
If it starts above 0, the differences between bars are visually exaggerated.
For example, in this activity, starting the axis at 30\% would make Netflix
look dominant when Disney+ and Hulu are also significant. A graph that starts
at 0 gives an honest picture of the data.}

\begin{teachernote}
Accept any of the three misleading tricks (non-zero axis, 3-D, inconsistent
scale), but the non-zero axis is the most commonly tested on AP exams.
Full credit requires: naming the feature, explaining the effect, and connecting
to a specific example.
\end{teachernote}
\end{tcolorbox}

\end{document}
